\section{Problem Formulation}

% \subsection{The Challenges}

% Research in multi-turn intent recognition has predominantly focused on modeling strategies. However, a critical challenge often overlooked is the scarcity of annotated, application-specific multi-turn datasets. Publicly available datasets typically encompass a limited number of intents, generally in the tens, which does not reflect the complexity encountered in real-world applications requiring hundreds of intents to accurately respond to user needs. The necessity for large-scale, annotated training datasets for effective model performance is evident, yet the dynamic nature of intent lists in ongoing business operations complicates dataset maintenance due to frequent updates and expansions of intent categories.

% Furthermore, the annotation process for multi-turn conversations is considerably more complex and error-prone than that for single-turn interactions. This complexity arises from the nuanced roles of context in conversations, including dependencies such as system responses influencing subsequent user queries. These factors contribute to potential semantic drifts over time and complicate the training data curation process, especially when considering the need for multi-lingual support.

% \subsection{Single-turn vs. Multi-turn Intent Recognition}
\subsection{Hierarchical Text Classification}
Hierarchical text classification (HTC) is a type of text classification where the classes or categories are organized in a hierarchy or tree structure. Instead of having a flat list of categories, the categories are arranged in a nested manner, so we need to consider the relationships of the nodes from different levels in the class taxonomy. 
% However, methods which optimally utilise the overall label hierarchy information often outperform the methods which simply disregard the structure and treat them as flat text classification, as they can take advantage of the extra information. \cite{Rojas2020EfficientSF}

The intents in our scenario are organised in a hierarchical tree structure. Specifically, each category can belong to at most one parent category and can have arbitrary number of children categories. Our class taxonomy $\mathcal{T}$ is a tree with fixed depth 3 and one meta root node, that is, the distance from the root node to all leaf nodes is 3, and the unique path from the root node to each leaf node will form one intent. 

We formulate the single-turn HTC task in our scenario as such, given a text $q_i$, the purpose is to predict a subset $I$ of the complete intent set $\mathcal{I}$, where the size of the subset $|I|$ containing the category from each layer is 3 (excluding root node). Generally, the number of intent sets exceeds 200.  

% \subsection{Single-turn Intent Classification} In the single-turn scenario, the objective is to classify a user's query $q$ into one of the predefined intent set $|I|$ from $\mathcal{I}$, where $|\mathcal{I}|$ often exceeds 200. Queries are complete and meaningful texts ranging from informational requests to action-oriented dialogues. Recognizing the correct intent enables the dialogue system to provide relevant solutions or perform specific actions. This process is further complicated in transnational applications that must accommodate queries in multiple languages and adapt to localized business operations.

\subsection{Multi-turn Intent Classification} Multi-turn scenario shares the same $\mathcal{T}$ and $\mathcal{I}$ as single-turn scenario. It involves a series of user queries $\mathcal{Q} = \{q_i\}_{i=1}^n$ in dialogues, the objective is to identify the intent of the final query $q_n$. Multi-turn recognition must account for the entire conversational context $\mathcal{C} = \{q_i\}_{i=1}^{n-1}$, which includes the historical queries. This context-dependency introduces additional complexity, requiring models to interpret nuanced conversational dynamics and adjust to evolving user intentions over the course of an interaction.

\subsection{Objective} This work aims to use easily accessible single-turn data to improve the accuracy of multi-turn intent recognition without requiring any multi-turn training datasets. 

% By proposing LARA that dynamically adapts to the conversation's context while minimizing reliance on large-scale annotated multi-turn datasets, we seek to mitigate the significant annotation challenges and enable more effective intent recognition in complex multi-turn dialogues.


% \subsection{Multi-turn Intent Recognition}
% While studies have been done on the multi-turn intent recognition, the focus is mainly on the modelling part. We argue that the most practical aspect of the task, which is the absence of annotated application-specific multi-turn dataset, is ignored. Moreover, the public datasets only have very few intents, usually on the scale of tens. When facing a real-life scenario where several hundreds of intents are needed to cater to users' need, a large scale annotated training dataset is needed for the supervised finetuned model to perform adequately good. One previous work \cite{Liu2020EnhancingMD} has considered the curation of large scale domain-specific multi-turn training data spanning a large number of intents. However, it requires manual annotation of 600k multi-turn sessions firstly to train a model to automatically annotate the remaining 400k sessions, making up to 1M samples in total. Considering that for an on-going business, the intent list will most likely change from time to time, e.g. some existing intents might be further broken down into more granular ones, and this will require re-annotation of part of the massive dataset.

% On the other hand, in single-turn intent recognition, the training data is easier to be constructed and maintained. This is due to there only exists one query and intent pair for each training sample. The complex behaviour and roles of context in a multi-turn conversation \cite{ghosal-etal-2021-exploring} makes the annotation process difficult and tricky. There are also other dependency possibilities to be considered for inclusion such as the influence of system answer on the next user query, i.e. the solutions shown are important for context understanding, but they might change over time, causing the possible user follow-up queries to change, and hence a semantic drift in the future. Not to mention about the dimension of query languages, which will need more annotated samples for training. In this work, we aim to find a relatively hassle-free method with strong performance which only extends from the available single-turn training data. 

% In the remaining subsections, we will formulate the problem for the ease of further discussions.

% \subsubsection{Single-turn Intent Recognition} In a single-turn scenario, given a user query $q$, the goal is to recognize the user intention from a list of intent classes $\mathcal{I} = \{I_i\}_{i=1}^k$, where $k$ is usually more than 200. A query could be an answer-seeking question like enquiries about the payment methods supported, or a task-oriented dialogue like order cancellation. By recognizing the intent, the dialogue system can help to provide the solution or perform the action accordingly. In a transnational business application, the model is expected to handle queries in multiple and mixed languages. The intent list $\mathcal{I}$ for each market could be also different to cater to the local business operations. In our scenario, each intent $I$ comprises $\langle r, l \rangle$ where $r$ is a representative query for the intent in the local language, and $l$ is the hierarchical label names in English (``\emph{Product} $>$ \emph{intention}", e.g. ``Food $>$ Cancel order", ``Digital Product $>$ Cancel order", ...).

% % \subsection{What’s single-turn/multi-turn intent recognition task}
% \subsubsection{Multi-turn Intent Recognition} In a multi-turn conversation scenario, we are given a conversation session corresponding to a sequence of queries, i.e., $\mathcal{Q} = \{q_i\}_{i=1}^n$. The goal is to recognize the user intention for the last query $q_n$. Each query turn can be ambiguous and is usually context-dependent, thus we need to consider the conversational context $\mathcal{C}$ when recognizing the user's true intention. Two types of context $\mathcal{C}$ are considered in the methods to be presented: (1) The context containing historical queries and the corresponding ground truth intent $\mathcal{C}_{\mathcal{Q}\mathcal{I}^*} = \{q_i, I^*_i\}_i^{n-1}$ and (2) containing historical queries only $\mathcal{C}_{\mathcal{Q}} = \{q_i\}_i^{n-1}$. During inference time, when the ground truth $I^*_i$ is not known, it is replace with the model prediction for that corresponding turn. 

% Both single-turn and multi-turn share the same intent classes $\mathcal{I}$. We assume that the training data $\mathcal{D} = \{\langle x_i, y_i\rangle\}$ for single-turn intent recognition, where $y_i \in \mathcal{I}$, is readily available as it is easier to be curated compared to its multi-turn counter part, and a traditional classification model $\mathcal{M}_c$ is trained on the dataset $\mathcal{D}$.
